\chapter{Point Processes and Hawkes Processes}\label{ch:qm:point-processes-and-hawkes-processes}

Count the trades in a liquid futures contract over a busy morning in one-minute bins. If trades
arrived like a Poisson clock, the variance of the counts would equal their mean; it is several
times larger, and the trades come in bursts, each trade making the next one likelier within a
second or so. A model in which every event raises the rate of further events, the \omterm{def:qm:point-processes-and-hawkes-processes:hawkes}{Hawkes
process}, reproduces that, and one number summarises it: the \omterm{def:qm:point-processes-and-hawkes-processes:hawkes}{branching ratio}, the average number
of trades each trade triggers. Fitted to the E-mini S\&P futures, it has been reported close to
one (Hardiman, Bercot and Bouchaud, 2013) and rising from about 0.3 to above 0.7 over a decade (Filimonov and
Sornette, 2012), and the same authors have shown how easily a fit invents excitation that is not
there. This chapter defines point processes by their intensities, builds the \omterm{def:qm:point-processes-and-hawkes-processes:hawkes}{Hawkes process}, fits
it by maximum likelihood in linear time, simulates it two ways, checks it by a change of clock, and
measures how a U-shaped intraday pattern fools the fit.

\section{Intensities and compensators}

\begin{definition}[Point process, counting process]\label{def:qm:point-processes-and-hawkes-processes:pp}
A \emph{point process}\index{point process} on $[0, \infty)$ is a random, locally finite set of event
times $t_1 < t_2 < \dots$. Its \emph{counting process}\index{counting process} is $N_t = \#\{i : t_i \le t\}$,
a right-continuous, integer-valued, nondecreasing \omterm{def:qm:probability-at-speed:filtration}{adapted process} with jumps of one.
\end{definition}

\begin{definition}[Conditional intensity, compensator]\label{def:qm:point-processes-and-hawkes-processes:intensity}
The \emph{conditional intensity}\index{conditional intensity} of a \omterm{def:qm:point-processes-and-hawkes-processes:pp}{counting process} is the \omterm{def:qm:probability-at-speed:filtration}{adapted process}
$\lambda_t$ with $\P(N_{t+dt} - N_t = 1 \mid \mathcal F_t) = \lambda_t\,dt + o(dt)$: the instantaneous rate of an event
given the past. Its \emph{compensator}\index{compensator} is $\Lambda_t = \int_0^t\lambda_s\,ds$, the predictable
increasing process with $N_t - \Lambda_t$ a \omterm{def:qm:ito-calculus:local}{local martingale}.
\end{definition}

The Doob--Meyer decomposition guarantees that every \omterm{def:qm:point-processes-and-hawkes-processes:pp}{counting process} has a \omterm{def:qm:point-processes-and-hawkes-processes:intensity}{compensator}; when it is
absolutely continuous, its density is the intensity. The intensity is the series' $\lambda_t$: the Poisson
rate of chapter~6, the hazard rate of a default time in One Quant Book~2, and the Hawkes intensity
below are all special cases.

\begin{definition}[Inhomogeneous Poisson process, Cox process]\label{def:qm:point-processes-and-hawkes-processes:cox}
An \emph{inhomogeneous Poisson process}\index{inhomogeneous Poisson process} has a deterministic intensity
$\lambda(t)$: counts in disjoint intervals are independent Poisson with means $\int\lambda(t)\,dt$. A \emph{Cox
process}\index{Cox process} (doubly stochastic \omterm{def:qm:jump-processes:poisson}{Poisson process}) is a \omterm{def:qm:jump-processes:poisson}{Poisson process} whose intensity is
itself a random process, independent of the events it drives.
\end{definition}

A \omterm{def:qm:point-processes-and-hawkes-processes:cox}{Cox process}'s counts are overdispersed: $\Var(N_t) = \E[\Lambda_t] + \Var(\Lambda_t)$. Book~6 models defaults
this way, with a random hazard rate. The randomness of a Cox intensity comes from outside; the Hawkes
intensity is driven by the process's own events.

\section{Self- and mutual excitation}

\begin{definition}[Hawkes process, excitation kernel, branching ratio]\label{def:qm:point-processes-and-hawkes-processes:hawkes}
A \emph{Hawkes process}\index{Hawkes process} is a \omterm{def:qm:point-processes-and-hawkes-processes:pp}{counting process} with intensity
\[
\lambda_t = \mu + \sum_{t_i < t}g(t - t_i),
\]
with a baseline $\mu > 0$ and a nonnegative \emph{excitation kernel}\index{excitation kernel} $g$. The
\emph{branching ratio}\index{branching ratio} is $\lVert g\rVert_1 = \int_0^\infty g(u)\,du$. With the exponential
kernel $g(u) = \alpha e^{-\beta u}$, it is $\alpha/\beta$.
\end{definition}

\begin{proposition}[Stationarity and the mean intensity]\label{prop:qm:point-processes-and-hawkes-processes:mean}
If $\lVert g\rVert_1 < 1$, a stationary \omterm{def:qm:point-processes-and-hawkes-processes:hawkes}{Hawkes process} exists and its mean intensity is $\bar\lambda =
\mu/(1 - \lVert g\rVert_1)$.
\end{proposition}

\begin{proof}
Taking expectations in the definition, with $\E[dN_s] = \bar\lambda\,ds$ in the stationary regime,
$\bar\lambda = \mu + \int_{-\infty}^t g(t - s)\bar\lambda\,ds = \mu + \bar\lambda\lVert g\rVert_1$. Existence is the
cluster construction below, whose clusters are finite when $\lVert g\rVert_1 < 1$.
\end{proof}

\begin{theorem}[Cluster representation]\label{thm:qm:point-processes-and-hawkes-processes:cluster}
A stationary \omterm{def:qm:point-processes-and-hawkes-processes:hawkes}{Hawkes process} is a Poisson cluster process: immigrants arrive as a \omterm{def:qm:jump-processes:poisson}{Poisson process} of rate
$\mu$; each event, immigrant or not, has a Poisson number of children with mean $\lVert g\rVert_1$, born at
independent delays with density $g/\lVert g\rVert_1$.
\end{theorem}

\admitted

Hawkes and Oakes (1974) proved it. Each immigrant starts a Galton--Watson tree whose expected size is
$1/(1 - \lVert g\rVert_1)$, so a share $\lVert g\rVert_1$ of all events are children: the \omterm{def:qm:point-processes-and-hawkes-processes:hawkes}{branching ratio} is the
fraction of activity that is endogenous (\cref{fig:qm:point-processes-and-hawkes-processes:cluster}). At
0.7 each exogenous trade brings 3.3 trades in all. The counts show it: for the exponential kernel, the
variance-to-mean ratio of counts in a window of length $\tau$ rises from one to $1/(1 - \lVert g\rVert_1)^2$ as the
window grows, from the covariance density $c(u) = \bar\lambda\frac{n(2-n)\beta}{2(1-n)}e^{-\beta(1-n)|u|}$ with $n =
\alpha/\beta$ (\cref{fig:qm:point-processes-and-hawkes-processes:dispersion}).

\begin{omfigure}
\begin{tikzpicture}[x=1.25cm,y=0.9cm,font=\small]
\draw[->,thick,gray] (0,0) -- (11,0) node[right,black] {time};
\foreach \x/\n in {0.8/a,4.3/b,8.2/c} {\draw[omDef,very thick] (\x,0) -- (\x,1.1); \fill[omDef] (\x,1.1) circle (2.5pt);}
\foreach \x/\p/\h in {1.4/0.8/0.8,1.7/0.8/0.8,2.3/1.7/0.55,2.6/1.4/0.55,3.1/2.6/0.35,4.9/4.3/0.8,5.8/4.9/0.55,8.5/8.2/0.8,8.9/8.2/0.8,9.4/8.9/0.55} {
  \draw[omThm,thick] (\x,0) -- (\x,\h); \fill[omThm] (\x,\h) circle (2pt);}
\draw[gray,-{Stealth}] (0.8,1.1) to[bend left=30] (1.4,0.85);
\draw[gray,-{Stealth}] (0.8,1.1) to[bend left=30] (1.7,0.85);
\draw[gray,-{Stealth}] (1.7,0.85) to[bend left=30] (2.3,0.6);
\draw[gray,-{Stealth}] (1.4,0.85) to[bend left=40] (2.6,0.6);
\draw[gray,-{Stealth}] (2.6,0.6) to[bend left=30] (3.1,0.4);
\draw[gray,-{Stealth}] (4.3,1.1) to[bend left=30] (4.9,0.85);
\draw[gray,-{Stealth}] (4.9,0.85) to[bend left=30] (5.8,0.6);
\draw[gray,-{Stealth}] (8.2,1.1) to[bend left=30] (8.5,0.85);
\draw[gray,-{Stealth}] (8.2,1.1) to[bend left=30] (8.9,0.85);
\draw[gray,-{Stealth}] (8.9,0.85) to[bend left=30] (9.4,0.6);
\node[omDef,font=\footnotesize,anchor=south] at (4.3,1.25) {immigrant (rate $\mu$)};
\node[omThm,font=\footnotesize,anchor=north] at (5.8,-0.15) {offspring: $\text{Poisson}(\lVert g\rVert_1)$ per event, delays $\propto g$};
\end{tikzpicture}

\omcaption{The cluster representation of a \omterm{def:qm:point-processes-and-hawkes-processes:hawkes}{Hawkes process}. Immigrants (tall, blue) arrive as a \omterm{def:qm:jump-processes:poisson}{Poisson
process}; every event has a Poisson number of children (short, red) with mean equal to the \omterm{def:qm:point-processes-and-hawkes-processes:hawkes}{branching ratio},
at delays drawn from the normalised kernel. The events of all generations together form the \omterm{def:qm:point-processes-and-hawkes-processes:hawkes}{Hawkes
process}.}\label{fig:qm:point-processes-and-hawkes-processes:cluster}
\end{omfigure}

\begin{definition}[Multivariate Hawkes process]\label{def:qm:point-processes-and-hawkes-processes:multi}
A \emph{multivariate Hawkes process}\index{multivariate Hawkes process} has components $N^{(1)}, \dots,
N^{(d)}$ with intensities $\lambda^{(i)}_t = \mu_i + \sum_j\sum_{t^{(j)}_k < t}g_{ij}(t - t^{(j)}_k)$; the branching matrix
is $G = (\lVert g_{ij}\rVert_1)$, and a stationary version exists when its spectral radius is below one.
\end{definition}

Buys exciting sells, trades exciting quote revisions, one venue's activity exciting another's: mutual
excitation is how the order-flow models of One Quant Books 10 and 11 are written.

\begin{example}[Buys and sells that excite each other]\label{ex:qm:point-processes-and-hawkes-processes:bivariate}
Let buys and sells each have baseline $0.1$ per second and exponential kernels with $\beta = 1$, a buy raising the
buy intensity by $0.2$ and the sell intensity by $0.4$, and symmetrically. The branching matrix $G = \bigl(\begin{smallmatrix}0.2 &
0.4\\0.4 & 0.2\end{smallmatrix}\bigr)$ has eigenvalues $0.6$ and $-0.2$, so the process is stable; taking expectations as in
\cref{prop:qm:point-processes-and-hawkes-processes:mean}, the mean intensities solve $\bar\lambda = \mu + G\bar\lambda$, giving
$\bar\lambda = (I_2 - G)^{-1}\mu = (0.25, 0.25)$ per second. Cross-excitation dominates: of the triggered activity on each side,
two thirds comes from the other side's trades.
\end{example}

\begin{omfigure}
\begin{tikzpicture}
\begin{groupplot}[group style={group size=2 by 1,horizontal sep=1.6cm},
  width=7.6cm,height=4.8cm,tick label style={font=\small},label style={font=\small},
  grid=major,grid style={gray!25},table/col sep=comma]
\nextgroupplot[xlabel={seconds},ylabel={$\lambda_t$ (per second)},xmin=0,xmax=60,ymin=-0.3,ymax=2.1]
\addplot[omDef,thick] table[x=s,y=lam]{figdata/methods/07-point-processes-and-hawkes-processes/intensity.csv};
\addplot[omThm,only marks,mark=|,mark size=3pt] table[x=s,y=y]{figdata/methods/07-point-processes-and-hawkes-processes/events.csv};
\addplot[gray,dashed,domain=0:60] {0.1};
\nextgroupplot[xlabel={bin length $\tau$ (seconds)},ylabel={$\Var N(\tau)/\E N(\tau)$},xmode=log,xmin=0.2,xmax=400,ymin=0,ymax=12,
  legend style={font=\footnotesize,at={(0.03,0.97)},anchor=north west,fill=white,draw=none}]
\addplot[omDef,very thick] table[x=tau,y=theory]{figdata/methods/07-point-processes-and-hawkes-processes/dispersion.csv};
\addplot[omThm,only marks,mark=*,mark size=1.6pt] table[x=tau,y=sim]{figdata/methods/07-point-processes-and-hawkes-processes/dispersion.csv};
\addplot[gray,dashed,domain=0.2:400] {1};
\legend{theory,simulated}
\end{groupplot}
\end{tikzpicture}

\omcaption{Left: one minute of a simulated trading day with $\mu = 0.1$, $\alpha = 0.7$ and $\beta = 1$ per second
(\omterm{def:qm:point-processes-and-hawkes-processes:hawkes}{branching ratio} 0.7): the intensity jumps by $\alpha$ at each event (ticks) and decays toward the baseline
(dashed). Right: variance-to-mean ratio of the counts in bins of length $\tau$, from the covariance density
and in the simulated day, against one for a \omterm{def:qm:jump-processes:poisson}{Poisson process} (dashed), rising to $1/(1 - 0.7)^2 = 11.1$. Data: the chapter's tutorial,
seeded.}\label{fig:qm:point-processes-and-hawkes-processes:dispersion}
\end{omfigure}

\section{Likelihood estimation}

\begin{proposition}[Log-likelihood of a point process]\label{prop:qm:point-processes-and-hawkes-processes:loglik}
For events $t_1 < \dots < t_n$ on $[0, T]$ and a parametric intensity $\lambda_t(\theta)$,
\[
\ell_n(\theta) = \sum_{i=1}^n\ln\lambda_{t_i}(\theta) - \int_0^T\lambda_t(\theta)\,dt.
\]
For the exponential \omterm{def:qm:point-processes-and-hawkes-processes:hawkes}{Hawkes process}, with $A_i = \sum_{j<i}e^{-\beta(t_i - t_j)}$,
\[
\ell_n = \sum_i\ln(\mu + \alpha A_i) - \mu T - \frac\alpha\beta\sum_i\bigl(1 - e^{-\beta(T - t_i)}\bigr), \qquad A_i =
e^{-\beta(t_i - t_{i-1})}(1 + A_{i-1}).
\]
\end{proposition}

\begin{proof}
The density of the events is the product, over the gaps, of the probability of no event, $e^{-\int\lambda}$,
times the rate at each event, $\lambda_{t_i}$: the likelihood of a sequence of exponential waiting times with a
time-varying rate. The recursion follows from $e^{-\beta(t_i - t_j)} = e^{-\beta(t_i - t_{i-1})}e^{-\beta(t_{i-1} - t_j)}$.
\end{proof}

The recursion makes the likelihood cost linear in the number of events, instead of quadratic; it is the
kernel of the running project's estimator, in Python, C++20 and Rust (the
tutorial's second listing). The maximum is found numerically: on one
simulated day of 8\,006 events with $n = 0.7$, the estimate is $(\hat\mu, \hat\alpha, \hat\beta) = (0.103, 0.696, 0.995)$,
a \omterm{def:qm:point-processes-and-hawkes-processes:hawkes}{branching ratio} of 0.700; over forty simulated days the estimates have mean 0.700 and standard deviation
0.011. Chapter~11 derives the standard errors that maximum likelihood carries.

\section{Simulation by thinning and by branching}

\begin{definition}[Thinning]\label{def:qm:point-processes-and-hawkes-processes:thinning}
\emph{Thinning}\index{thinning} simulates a \omterm{def:qm:point-processes-and-hawkes-processes:pp}{point process} with intensity $\lambda_t \le M_t$, where $M$ is a bound known
in advance, by proposing events at rate $M$ and accepting a proposal at $t$ with probability $\lambda_t/M_t$.
\end{definition}

\begin{method}[Ogata's thinning for a Hawkes process]\label{met:qm:point-processes-and-hawkes-processes:ogata}
Between events the exponential Hawkes intensity decreases, so its value just after the last event bounds it
until the next. Draw a waiting time at that rate, decay the excitation over it, accept with probability
$\lambda/\text{bound}$, and on acceptance add $\alpha$ to the excitation. The simulation is exact, with no time step (the
tutorial's first listing). Alternatively, simulate the clusters of
\cref{thm:qm:point-processes-and-hawkes-processes:cluster} generation by generation: the share of children
among the events is then observable, 0.699 on a simulated day against 0.7.
\end{method}

\section{Diagnostics by time change}

\begin{theorem}[Time-rescaling]\label{thm:qm:point-processes-and-hawkes-processes:rescaling}
If $N$ has continuous \omterm{def:qm:point-processes-and-hawkes-processes:intensity}{compensator} $\Lambda$ with $\Lambda_\infty = \infty$, the rescaled times $\Lambda(t_1) < \Lambda(t_2) < \dots$
form a \omterm{def:qm:jump-processes:poisson}{Poisson process} of unit rate: the gaps $\Lambda(t_i) - \Lambda(t_{i-1})$ are independent $\mathrm{Exp}(1)$.
\end{theorem}

\admitted

Running the clock at the speed of the model's intensity turns any correctly specified \omterm{def:qm:point-processes-and-hawkes-processes:pp}{point process} into a
standard \omterm{def:qm:jump-processes:poisson}{Poisson process}, so the fitted \omterm{def:qm:point-processes-and-hawkes-processes:intensity}{compensator} gives residuals to test: their mean should be one, their
variance one, their quantiles those of the exponential law (Brown et al., 2002). For the fitted Hawkes model
the residuals have mean 1.000 and variance 0.98; for a Poisson model with the same mean rate, the variance is
4.63 and the upper quantiles are far too large (\cref{fig:qm:point-processes-and-hawkes-processes:qq}).

\begin{omfigure}
\begin{tikzpicture}
\begin{axis}[width=8.5cm,height=6cm,
  xlabel={$\mathrm{Exp}(1)$ quantile},ylabel={residual quantile},
  tick label style={font=\small},label style={font=\small},
  xmin=0,xmax=10,ymin=0,ymax=28,grid=major,grid style={gray!25},
  legend style={font=\footnotesize,at={(0.03,0.97)},anchor=north west,fill=white,draw=none},
  table/col sep=comma]
\addplot[omThm,thick,mark=*,mark size=1.1pt] table[x=q,y=poisson]{figdata/methods/07-point-processes-and-hawkes-processes/qq.csv};
\addplot[omDef,thick,mark=*,mark size=1.1pt] table[x=q,y=hawkes]{figdata/methods/07-point-processes-and-hawkes-processes/qq.csv};
\addplot[gray,dashed,domain=0:10] {x};
\legend{Poisson fit,Hawkes fit,45-degree line}
\end{axis}
\end{tikzpicture}

\omcaption{Quantiles of the time-rescaled residuals of one simulated day against those of the unit
exponential. The fitted Hawkes model lies on the diagonal; a Poisson model with the same mean rate misses the
bursts, and its largest gaps are almost three times too long. Data: the chapter's tutorial,
seeded.}\label{fig:qm:point-processes-and-hawkes-processes:qq}
\end{omfigure}

\section{Tutorial: the self-exciting tape}\label{tut:qm:point-processes-and-hawkes-processes:tape}

\begin{tutorial}
\textbf{Goal.} Simulate a day of self-exciting trades, fit it by maximum likelihood, check the fit by time
rescaling, and fit a day with no excitation but an intraday pattern. \textbf{End state:}
\cref{fig:qm:point-processes-and-hawkes-processes:dispersion,fig:qm:point-processes-and-hawkes-processes:qq}
and the \omterm{def:qm:point-processes-and-hawkes-processes:hawkes}{branching ratios} of the weekend problem.
\begin{tutsteps}
\item \textbf{Simulate} by Ogata's \omterm{def:qm:point-processes-and-hawkes-processes:thinning}{thinning}.
\omcode{code/firm/hawkes/firm_hawkes.py}{24}{40}{Ogata's thinning for the exponential Hawkes process.}\label{lst:qm:point-processes-and-hawkes-processes:thinning}
\item \textbf{Fit}: the log-likelihood in linear time, maximised by a Nelder--Mead search over $(\ln\mu,
\operatorname{logit}n, \ln\beta)$ so that the \omterm{def:qm:point-processes-and-hawkes-processes:hawkes}{branching ratio} stays below one.
\omcode{code/firm/hawkes/firm_hawkes.py}{98}{111}{The log-likelihood with the recursion for $A_i$.}\label{lst:qm:point-processes-and-hawkes-processes:loglik}
\item \textbf{Run} \texttt{problem()} and \texttt{fig\_hawkes.py}; the C++20 and Rust twins in
\texttt{code/firm/hawkes/} reproduce the Python log-likelihood to $10^{-12}$ on a fixed event list.
\end{tutsteps}
\textbf{What to change next.} Fit a two-exponential kernel and see whether the fast and slow parts separate;
simulate a bivariate process of buys and sells that excite each other and recover the branching matrix.
\end{tutorial}

\section{Build: the Hawkes toolkit}\label{bld:qm:point-processes-and-hawkes-processes:hawkes}

\begin{build}
\bfield{Purpose} Model the arrival of trades, orders and cancellations in the miniature firm's markets:
simulate realistic bursty flow for the exchange simulator, and estimate excitation from recorded
flow.
\bfield{Interface} Python: \texttt{simulate\_thinning}, \texttt{simulate\_branching},
\texttt{simulate\_multivariate}, \texttt{intensity}, \texttt{loglik}, \texttt{fit}, \texttt{compensator},
\texttt{residuals}, \texttt{nelder\_mead}. C++20 and Rust: \texttt{Params}, \texttt{simulate\_thinning},
\texttt{loglik}, \texttt{compensator}.
\bfield{Rules} Exact simulation (no time step); likelihood in $O(n)$; parameters searched on a scale that
keeps $\mu, \beta > 0$ and $0 \le n < 1$; the three languages agree to $10^{-12}$ on the reference list.
\bfield{Acceptance tests} \texttt{code/firm/hawkes/\{tests,cpp,rust\}}: mean counts equal $\mu T/(1 - n)$;
maximum likelihood recovers the parameters; residuals of the true model are unit exponential; branching and
\omterm{def:qm:point-processes-and-hawkes-processes:thinning}{thinning} simulations have the same count distribution.
\bfield{Stretch} Multivariate maximum likelihood; power-law kernels; a time-varying baseline $\mu(t)$ estimated
jointly with the kernel (the weekend problem's fix).
\end{build}

\begin{omsources}
\item A.~G.~Hawkes, ``Spectra of some self-exciting and mutually exciting point processes'',
\emph{Biometrika} 58, 1971.
\item A.~G.~Hawkes and D.~Oakes, ``A cluster process representation of a self-exciting process'',
\emph{Journal of Applied Probability} 11, 1974.
\item Y.~Ogata, ``On Lewis' simulation method for point processes'', \emph{IEEE Transactions on
Information Theory} 27, 1981.
\item E.~N.~Brown, R.~Barbieri, V.~Ventura, R.~E.~Kass and L.~M.~Frank, ``The time-rescaling theorem and
its application to neural spike train data analysis'', \emph{Neural Computation} 14, 2002.
\item V.~Filimonov and D.~Sornette, ``Quantifying reflexivity in financial markets'', \emph{Physical Review E}
85, 2012; ``Apparent criticality and calibration issues in the Hawkes self-excited point process model'',
\emph{Quantitative Finance} 15, 2015.
\item S.~J.~Hardiman, N.~Bercot and J.-P.~Bouchaud, ``Critical reflexivity in financial markets: a Hawkes
process analysis'', \emph{European Physical Journal B} 86, 2013.
\end{omsources}

\section{Exercises}

\begin{exercise}[$\star$]\label{exo:qm:point-processes-and-hawkes-processes:1}
A \omterm{def:qm:point-processes-and-hawkes-processes:hawkes}{Hawkes process} has $\mu = 0.1$ per second, $\alpha = 0.7$ and $\beta = 1$ per second. What are its \omterm{def:qm:point-processes-and-hawkes-processes:hawkes}{branching ratio}, mean
intensity and expected number of events in a 23\,400-second trading day?
\end{exercise}

\begin{exercise}[$\star$]\label{exo:qm:point-processes-and-hawkes-processes:2}
In the same process, how many events does one immigrant bring on average, and what share of events are
immigrants?
\end{exercise}

\begin{exercise}[$\star$]\label{exo:qm:point-processes-and-hawkes-processes:3}
An \omterm{def:qm:point-processes-and-hawkes-processes:cox}{inhomogeneous Poisson process} has intensity $\lambda(t) = a + bt$. Write its \omterm{def:qm:point-processes-and-hawkes-processes:intensity}{compensator}, and the probability of
no event in $[0, T]$.
\end{exercise}

\begin{exercise}[$\star\star$]\label{exo:qm:point-processes-and-hawkes-processes:4}
How long after an event does its contribution to the intensity take to halve, with $\beta = 1$ per second? How large is
the jump of the intensity at an event compared with the baseline?
\end{exercise}

\begin{exercise}[$\star\star$]\label{exo:qm:point-processes-and-hawkes-processes:5}
Compute the variance-to-mean ratio of counts in bins of 1 and 60 seconds, and its limit, for the process of
\cref{exo:qm:point-processes-and-hawkes-processes:1}.
\end{exercise}

\begin{exercise}[$\star\star$]\label{exo:qm:point-processes-and-hawkes-processes:6}
Each day the arrival rate is 1 or 3 per second with equal probability, and within a day arrivals are Poisson.
What are the mean and the variance of the count in one second? Is this a \omterm{def:qm:point-processes-and-hawkes-processes:hawkes}{Hawkes process}?
\end{exercise}

\begin{exercise}[$\star\star\star$]\label{exo:qm:point-processes-and-hawkes-processes:7}
\emph{Coding.} Fit the exponential Hawkes model to forty simulated days of the process of
\cref{exo:qm:point-processes-and-hawkes-processes:1} and report the mean and standard deviation of the
estimated \omterm{def:qm:point-processes-and-hawkes-processes:hawkes}{branching ratio}.
\end{exercise}

\begin{exercise}[$\star\star\star$]\label{exo:qm:point-processes-and-hawkes-processes:8}
\emph{Find the flaw.} ``We fitted a Hawkes model to a full day of trades and found a \omterm{def:qm:point-processes-and-hawkes-processes:hawkes}{branching ratio} of 0.87:
the market is close to critical.''
\end{exercise}

\section{Problem: The Self-Exciting Tape}

\begin{problem}[{Weekend problem --- excitation that is there, and excitation that is not}]\label{pb:qm:point-processes-and-hawkes-processes:1}
Trades in a contract arrive over a 23\,400-second day as an exponential \omterm{def:qm:point-processes-and-hawkes-processes:hawkes}{Hawkes process} with $\mu = 0.1$, $\alpha =
0.7$ and $\beta = 1$ per second. A second day has no excitation at all: trades are an \omterm{def:qm:point-processes-and-hawkes-processes:cox}{inhomogeneous Poisson process}
whose rate is three times higher at the open and the close than at midday, $\lambda(t) \propto 1 + 8(t/T -
\tfrac12)^2$, with the same expected number of trades.

\medskip\noindent\textbf{Part I --- The model.}
\begin{enumerate}
\item What is the \omterm{def:qm:point-processes-and-hawkes-processes:hawkes}{branching ratio}?
\item What are the mean intensity and the expected number of trades in the day?
\item What is the expected cluster size, and the share of immigrants?
\item What is the \omterm{def:qm:stochastic-differential-equations:ou}{half-life} of an event's excitation?
\item What are the variance-to-mean ratios of one-second and one-minute counts?
\end{enumerate}

\medskip\noindent\textbf{Part II --- Estimation.}
\begin{enumerate}[resume]
\item How many trades does the simulated day have?
\item What does maximum likelihood estimate?
\item Over forty simulated days, what are the mean and standard deviation of the estimated \omterm{def:qm:point-processes-and-hawkes-processes:hawkes}{branching ratio}?
\item What mean and variance do the time-rescaled residuals have, for the Hawkes fit and for a Poisson fit?
\item In a branching simulation of the day, what share of events are children?
\end{enumerate}

\medskip\noindent\textbf{Part III --- The seasonal day.}
\begin{enumerate}[resume]
\item What is the rate at the open relative to midday, and the mean rate?
\item How many trades does the seasonal day have?
\item What \omterm{def:qm:point-processes-and-hawkes-processes:hawkes}{branching ratio} and decay rate does the Hawkes fit find?
\item Over twenty seasonal days, and over twenty days of a constant-rate \omterm{def:qm:jump-processes:poisson}{Poisson process}?
\item Why does the fit find excitation where there is none?
\end{enumerate}

\medskip\noindent\textbf{Part IV --- Judgement.}
\begin{enumerate}[resume]
\item How would you remove the bias?
\item What do the conflicting published estimates for the E-mini suggest?
\item How could a trading desk use a well-fitted Hawkes model?
\item State the \emph{named result}: the \omterm{def:qm:point-processes-and-hawkes-processes:hawkes}{branching ratio} recovered, and the spurious one.
\item In one sentence: when is a high \omterm{def:qm:point-processes-and-hawkes-processes:hawkes}{branching ratio} evidence of self-excitation?
\end{enumerate}
\end{problem}

\section{Interview questions}

\begin{interviewq}[$\star$ \iqroles{researcher, trader}]\label{iq:qm:point-processes-and-hawkes-processes:1}
What is a \omterm{def:qm:point-processes-and-hawkes-processes:hawkes}{Hawkes process}, and what does its \omterm{def:qm:point-processes-and-hawkes-processes:hawkes}{branching ratio} measure?
\end{interviewq}

\begin{interviewq}[$\star\star$ \iqroles{researcher}]\label{iq:qm:point-processes-and-hawkes-processes:2}
Show that the mean intensity of a stationary \omterm{def:qm:point-processes-and-hawkes-processes:hawkes}{Hawkes process} is $\mu/(1 - n)$. What happens when $n \ge 1$?
\end{interviewq}

\begin{interviewq}[$\star\star$ \iqroles{researcher, developer}]\label{iq:qm:point-processes-and-hawkes-processes:3}
How do you simulate a \omterm{def:qm:point-processes-and-hawkes-processes:hawkes}{Hawkes process} exactly?
\end{interviewq}

\begin{interviewq}[$\star\star$ \iqroles{developer}]\label{iq:qm:point-processes-and-hawkes-processes:4}
The Hawkes log-likelihood looks quadratic in the number of events. How do you evaluate it in linear time for an
exponential kernel?
\end{interviewq}

\begin{interviewq}[$\star\star$ \iqroles{trader, researcher}]\label{iq:qm:point-processes-and-hawkes-processes:5}
How would a market maker use an estimate of trade self-excitation?
\end{interviewq}

\begin{interviewq}[$\star\star\star$ \iqroles{researcher, mle}]\label{iq:qm:point-processes-and-hawkes-processes:6}
How would you test whether a fitted point-process model is adequate?
\end{interviewq}
